%% kapa: exam=1 round=27 year=2016 period=2 form=A subject=accounting subject_ko=회계학 seq=11 number=51
%% source: Q-Net 감정평가사 기출문제 (cst00309) artlSeq=5169402 file=extracted/1008/5169402/제27회 감정평가사 1차 문제/제27회 감정평가사 2교시문제(A형).pdf
%% conversion: kordoc 4.24.51 -> tools/gen1.py (auto); review=proofread-sonnet; flags=none
\begin{question}{51}
재무제표 표시에 관한 설명으로 옳지 않은 것은?
\begin{choices}
  \item 계속기업의 가정이 적절한지의 여부를 평가할 때 경영진은 적어도 보고기간말로부터 향후 12개월 기간에 대하여 이용가능한 모든 정보를 고려한다.
  \item 기업이 재무상태표에 유동자산과 비유동자산, 그리고 유동부채와 비유동부채로 구분하여 표시하는 경우, 이연법인세자산(부채)은 유동자산(부채)으로 분류하지 아니한다.
  \item 매입채무 그리고 종업원 및 그 밖의 영업원가에 대한 미지급비용과 같은 유동부채는 기업의 정상영업주기 내에 사용되는 운전자본의 일부이다. 이러한 항목은 보고기간 후 12개월 후에 결제일이 도래한다 하더라도 유동부채로 분류한다.
  \item 보고기간 후 12개월 이내에 만기가 도래하는 경우에는, 기업이 기존의 대출계약 조건에 따라 보고기간 후 적어도 12개월 이상 부채를 차환하거나 연장할 것으로 기대하고 있고, 그런 재량권이 있다고 하더라도, 유동부채로 분류한다.
  \item 비용을 기능별로 분류하는 기업은 감가상각비, 기타 상각비와 종업원급여비용을 포함하여 비용의 성격에 대한 추가 정보를 공시한다.
\end{choices}
\end{question}
\answer{4}
